\section{Agent Loop Implementations}
\label{app:loops}

We provide detailed pseudocode for each agent architecture. The actual implementations handle additional complexity including error recovery, optional tool calls (calculator, code interpreter), and provider-specific message formatting.

\subsection{\textsc{React} (Baseline)}

The baseline architecture follows the standard ReAct pattern: observe, reason, act. Tool results include the next observation appended to maintain conversational flow.

\begin{verbatim}
def run_episode(env, llm, optional_tools=[]):
    obs = env.reset()
    messages = [{"role": "user", "content": format_observation(obs, initial=True)}]
    all_tools = [ACTION_TOOL] + optional_tools
    
    while not obs.done:
        # Generate response - may use optional tools before final action
        while True:
            response = llm.generate(messages, tools=all_tools)
            messages.append(format_assistant_message(response))
            
            if response.tool_name == "take_action":
                break  # Final action submitted
            
            # Handle optional tool call (calculator, code interpreter)
            result = execute_tool(response.tool_name, response.tool_input)
            messages.append(format_tool_result(response.tool_id, result))
        
        # Execute action in environment
        action = parse_action(response.tool_input)
        obs = env.step(action)
        
        # Append tool result with next observation
        result_text = f"Sold {obs.cups_sold} cups, profit: ${obs.daily_profit/100:.2f}"
        tool_result = format_tool_result(response.tool_id, result_text)
        if not obs.done:
            tool_result["content"].append({"type": "text", 
                                           "text": format_observation(obs)})
        messages.append(tool_result)
    
    return obs.total_profit
\end{verbatim}

\subsection{\textsc{Plan-Act}}

Before each action, the agent generates an explicit multi-day plan. The planning phase uses no tools (text generation only); the action phase enables tool use.

\begin{verbatim}
PLANNING_PROMPT = """
Before taking action, create a strategic plan for the next 3-5 days:

1. **Weather Outlook**: What weather do you expect? How will it affect demand?
2. **Inventory Needs**: What supplies will you need? When should you buy them?
3. **Pricing Strategy**: What prices will work best for each weather condition?
4. **Location Considerations**: Should you consider changing location?
5. **Key Risks**: What could go wrong? How will you handle it?

Output your plan, then decide today's action based on it.
"""

def run_episode(env, llm, optional_tools=[]):
    obs = env.reset()
    messages = []
    all_tools = [ACTION_TOOL] + optional_tools
    
    while not obs.done:
        # Phase 1: Planning (no tools - text generation only)
        plan_prompt = format_observation(obs) + "\n\n" + PLANNING_PROMPT
        messages.append({"role": "user", "content": plan_prompt})
        
        plan_response = llm.generate(messages, tools=[])  # No tools
        plan_text = plan_response.text_content
        messages.append({"role": "assistant", "content": plan_text})
        
        # Phase 2: Action (with tools enabled)
        action_prompt = "Based on your plan above, now take today's action."
        messages.append({"role": "user", "content": action_prompt})
        
        response = generate_with_optional_tools(llm, messages, all_tools)
        action = parse_action(response.tool_input)
        obs = env.step(action)
        
        # Add result (next obs will be in next planning prompt)
        messages.append(format_tool_result(response.tool_id, 
                                           format_result(obs)))
    
    return obs.total_profit
\end{verbatim}

\subsection{\textsc{Act-Reflect}}

Inspired by Reflexion~\citep{shinn2023reflexion}, this architecture adds a reflection phase \emph{after} each action. The agent analyzes what worked and what could improve before seeing the next day's conditions.

\begin{verbatim}
REFLECTION_PROMPT = """
Reflect on yesterday's results before making today's decision:

1. **What Worked Well?** What decisions led to positive outcomes?
2. **What Could You Have Done Better?** Where did you lose money or 
   miss opportunities?
3. **What Surprised You?** Was demand higher or lower than expected?
4. **What Will You Do Differently?** How will you adjust your strategy?

Use these insights to inform today's action.
"""

def run_episode(env, llm, optional_tools=[]):
    obs = env.reset()
    messages = [{"role": "user", "content": format_observation(obs, initial=True)}]
    all_tools = [ACTION_TOOL] + optional_tools
    
    while not obs.done:
        # Action phase
        response = generate_with_optional_tools(llm, messages, all_tools)
        action = parse_action(response.tool_input)
        pre_obs = obs
        obs = env.step(action)
        
        # Add tool result
        result_text = format_result(obs)
        messages.append(format_tool_result(response.tool_id, result_text))
        
        if not obs.done:
            # Reflection phase (no tools - analysis only)
            reflect_prompt = f"""## Yesterday's Outcome
{result_text}

## Current State After Yesterday
- Cash: ${obs.cash/100:.2f}
- Total Profit: ${obs.total_profit/100:.2f}

{REFLECTION_PROMPT}"""
            messages.append({"role": "user", "content": reflect_prompt})
            
            reflection = llm.generate(messages, tools=[])  # No tools
            messages.append({"role": "assistant", "content": reflection.text})
            
            # Add next observation
            messages.append({"role": "user", 
                            "content": format_observation(obs)})
    
    return obs.total_profit
\end{verbatim}

\subsection{\textsc{Full} (Plan + Reflect)}

Combines both planning and reflection, following the pattern: Reflect (on previous day) $\rightarrow$ Plan (for upcoming days) $\rightarrow$ Act. This requires three LLM calls per turn.

\begin{verbatim}
PLAN_UPDATE_PROMPT = """
Review your previous plan in light of yesterday's results:

Previous plan: {previous_plan}

Has anything changed that invalidates this plan? Do you need to adjust 
your strategy? Update your plan if needed, then take today's action.
"""

def run_episode(env, llm, optional_tools=[]):
    obs = env.reset()
    messages = []
    current_plan = None
    all_tools = [ACTION_TOOL] + optional_tools
    
    while not obs.done:
        # Phase 1: Planning (includes observation)
        if current_plan:
            plan_prompt = format_observation(obs) + "\n\n" + \
                         PLAN_UPDATE_PROMPT.format(previous_plan=current_plan)
        else:
            plan_prompt = format_observation(obs) + "\n\n" + PLANNING_PROMPT
        
        messages.append({"role": "user", "content": plan_prompt})
        plan_response = llm.generate(messages, tools=[])
        current_plan = plan_response.text_content
        messages.append({"role": "assistant", "content": current_plan})
        
        # Phase 2: Action
        messages.append({"role": "user", 
                        "content": "Based on your plan, now take today's action."})
        response = generate_with_optional_tools(llm, messages, all_tools)
        action = parse_action(response.tool_input)
        obs = env.step(action)
        
        result_text = format_result(obs)
        messages.append(format_tool_result(response.tool_id, result_text))
        
        # Phase 3: Reflection (if game continues)
        if not obs.done:
            reflect_prompt = f"""## Yesterday's Outcome
{result_text}

## Your Previous Plan
{current_plan}

{REFLECTION_PROMPT}

Consider whether your plan is still valid or needs adjustment."""
            messages.append({"role": "user", "content": reflect_prompt})
            
            reflection = llm.generate(messages, tools=[])
            messages.append({"role": "assistant", "content": reflection.text})
    
    return obs.total_profit
\end{verbatim}

\subsection{Optional Tool Definitions}

Agents can optionally be provided with cognitive scaffolding tools. When enabled, agents may call these tools multiple times before submitting their final action.

\subsubsection{Calculator Tool}

Safe arithmetic evaluation supporting basic operations (+, -, *, /, //, \%, **):

\begin{verbatim}
{
  "name": "calculator",
  "description": "Perform arithmetic calculations. Useful for computing 
                  profit margins, break-even points, inventory costs, and 
                  other business math. Supports +, -, *, /, //, %, and ** 
                  operators.",
  "input_schema": {
    "type": "object",
    "properties": {
      "expression": {
        "type": "string",
        "description": "Mathematical expression to evaluate. Examples: 
                        '100 * 0.75', '(500 - 200) / 500 * 100', 
                        '12 * 4 * 0.25'"
      }
    },
    "required": ["expression"]
  }
}
\end{verbatim}

Implementation uses AST-based safe evaluation with guards against code injection, division by zero, and excessively large numbers.

\subsubsection{Code Interpreter Tool}

Sandboxed Python execution with access to mathematical libraries:

\begin{verbatim}
{
  "name": "run_python",
  "description": "Execute Python code for complex calculations, optimization, 
                  or data analysis. Has access to math, numpy (as np), and 
                  statistics modules. Use for profit calculations, demand 
                  forecasting, or break-even analysis.",
  "input_schema": {
    "type": "object",
    "properties": {
      "code": {
        "type": "string",
        "description": "Python code to execute. Has access to: math, numpy 
                        (np), statistics. Print results to see them. 
                        Example: print(sum([p * d for p, d in zip(prices, 
                        demands)]))"
      }
    },
    "required": ["code"]
  }
}
\end{verbatim}

Security measures include:
\begin{itemize}
    \item AST validation blocking dangerous operations (eval, exec, file I/O)
    \item Restricted builtins (no \texttt{type}, \texttt{getattr}, \texttt{\_\_import\_\_})
    \item 5-second execution timeout
    \item No network or filesystem access
    \item Blocked dunder attribute access (\texttt{\_\_class\_\_}, \texttt{\_\_bases\_\_}, etc.)
\end{itemize}

\subsection{Error Handling}

All architectures implement error recovery for invalid actions:

\begin{verbatim}
MAX_RETRIES = 3

def execute_with_retry(env, llm, messages, action):
    for retry in range(MAX_RETRIES):
        obs = env.step(action)
        
        if not obs.is_error_response:
            return obs  # Valid action
        
        # Feed error back to LLM
        error_msg = f"ACTION ERROR: {'; '.join(obs.action_errors)}"
        messages.append(format_error_feedback(error_msg, obs))
        
        # Get corrected action
        response = llm.generate(messages, tools=[ACTION_TOOL])
        action = parse_action(response.tool_input)
    
    # Fallback to safe default action
    return env.step(LemonadeAction(price_per_cup=50))
\end{verbatim}

This ensures agents can recover from budget constraint violations, invalid purchases, or malformed action parameters.

